% answers_short.tex —— 2010 年·普通高等学校招生考试·天津卷·理科 · 简短答案速查
% 由 tools/make_short_answers.py 自 answers.tex 的 \gksolhead/\gkanswer 提取，
% 并与 2010/2010_数学真题_天津理_答案要点.md 逐题比对后生成，勿手改（改源 MD 后重跑）。

\begin{gknotice}
\textbf{使用说明}：本节只给最终答案便于快速核对，完整推导见解析版。答案为编写者独立推导，\textbf{非官方标准答案}。
\end{gknotice}

\gkpart{选择题}
\begin{center}\begin{tabular}{ccccccccccc}
\toprule
题号 & 1 & 2 & 3 & 4 & 5 & 6 & 7 & 8 & 9 & 10 \\
\midrule
答案 & A & B & B & D & B & C & A & C & D & B \\
\bottomrule
\end{tabular}\end{center}

\gkpart{填空题}
\begin{itemize}[leftmargin=2.6em]
  \item[\textbf{11.}] $24$ 和 $23$。
  \item[\textbf{12.}] $\dfrac{10}3$。
  \item[\textbf{13.}] $(x+1)^2+y^2=2$。
  \item[\textbf{14.}] $\dfrac{\sqrt6}6$。
  \item[\textbf{15.}] $\sqrt3$。
  \item[\textbf{16.}] $\left(-\infty,-\dfrac{\sqrt3}2\right]\cup\left[\dfrac{\sqrt3}2,+\infty\right)$。
\end{itemize}

\gkpart{解答题}
\begin{itemize}[leftmargin=2.6em]
  \item[\textbf{17.}] （1）$T=\pi$，最大值 $2$，最小值 $-1$；（2）$\dfrac{3-4\sqrt3}{10}$。
  \item[\textbf{18.}] （1）$\dfrac{40}{243}$；（2）$\dfrac8{81}$；（3）$P(\xi=0)=\dfrac1{27}$，$P(\xi=1)=\dfrac29$，$P(\xi=2)=\dfrac4{27}$，$P(\xi=3)=\dfrac8{27}$，$P(\xi=6)=\dfrac8{27}$。
  \item[\textbf{19.}] （1）$\dfrac35$；（2）证明；（3）$\dfrac{\sqrt5}3$。
  \item[\textbf{20.}] （1）$\dfrac{x^2}4+y^2=1$；（2）$\pm\dfrac{2\sqrt{14}}5$。
  \item[\textbf{21.}] （1）$f$ 在 $(-\infty,1)$ 上单调递增，在 $(1,+\infty)$ 上单调递减，极大值 $f(1)=\dfrac1{\mathrm e}$，无极小值；（2）证明；（3）证明。
  \item[\textbf{22.}] （1）证明；（2）① 证明；② 证明。
\end{itemize}
